\section{Limitations}
\label{sec:limitations}
\begin{itemize}
\item \textbf{Ground-truth crops.} All benchmark results except the SPEED+ detector row use ground-truth boxes. On SPEED+, a synthetic-only detector more than doubles the rotation error on lightbox ($11.21^\circ\to23.88^\circ$) and raises the score from $0.227$ to $0.523$, mainly through poor real-domain boxes and the box-relative translation target (\cref{tab:spp_sota}). A deployable system therefore needs a detector that transfers to real images, or translation that does not depend on the box. Detection was not evaluated on the general-object benchmarks.
\item \textbf{Ground-truth translation in rotation inference.} On T-LESS and YCB-V the allocentric rotation is converted, and orbit candidates are mapped between views, with the ground-truth translation. The pre-registered evaluation with predicted translation, which also gives joint 6D success, has not yet run, so these rotation results are an upper reference for a system without ground-truth information.
\item \textbf{How general the shared architecture is.} The backbones, feature interface, head architecture, recipe and inference algorithm are shared, but each object needs its own labelled training images, intrinsics, crops and symmetry set, and its own heads; heads do not transfer zero-shot to new objects (\cref{sec:supp_zeroshot}). Two options are chosen per application on validation data: orbit aggregation helps on T-LESS, SYMSOL-T and SPEED+ but hurts on YCB-V, and MoGe-2 fusion helps only on SPEED+ and T-LESS. YCB-V also needed synthetic training data. Four applications are evidence for, not proof of, the claim that frozen backbones with small heads serve many applications; the industrial, warehouse and robotic settings named in the introduction were not evaluated.
\item \textbf{Matched comparisons.} Published numbers use different crops, masks, object sets and fine-tuned backbones; \cref{tab:tless_sota,tab:symsol_sota,tab:spp_sota,tab:trans_xobj} state each setting, and only rows with the same setting are matched.
\item \textbf{Hierarchy vs.\ grid.} The flat-baseline comparison of \cref{sec:abl_flat} uses a coarser grid than the tree; the same-grid control is still running, so the share of the gain due to hierarchical prediction is not yet established.
\item \textbf{Quantisation floor.} The depth-5 leaf quantisation ($2.51^\circ$ mean) bounds fine accuracy for asymmetric objects. Our best synthetic SPEED+ errors ($3.3$--$3.7^\circ$) are close to it, and it caps fine-threshold recall on T-LESS. A continuous refinement on top of the selected mode is a natural next step; the refiners we tried helped only on synthetic data (\cref{sec:supp_negative}).
\item \textbf{Cost.} The heads are small, but the system is not: the two backbones hold 634\,M parameters, feature extraction is the dominant adaptation cost, orbit aggregation multiplies inference cost by about 4.5, and the geometry branch by about 4. Head-only latency and reduced-precision accuracy are not yet measured. Timings come from a shared GPU and are upper bounds.
\item \textbf{Distillation} helps on SPEED+, but on SYMSOL-T it improves the single view only slightly (validation $2.54^\circ\to2.52^\circ$) and makes the orbit system worse ($2.07^\circ\to2.16^\circ$); its T-LESS evaluation is running.
\end{itemize}

\section{Conclusion}
\label{sec:conclusion}
Shared frozen appearance and geometry backbones can support 6D pose estimation of very different known objects through compact per-object heads. The same architecture, trained only in its 2.1--2.4\,M-parameter rotation trees and 0.33\,M-parameter translation block, reaches a recall of 79\% at $10^\circ$ and 88\% at 5\,cm on real T-LESS images, $2.07^\circ$ on SYMSOL-T, $12.9^\circ$ and 1.6\,cm on YCB-V, and $11.2^\circ/11.1^\circ$ on real spacecraft images after training on renders only. The decisive problem for such heads is selecting the right mode among symmetric or look-alike rotations. Multi-path targets address it for exact symmetries on every benchmark at no inference cost, and orbit aggregation and distillation address look-alike flips where aggregation helps, at a measured $4\text{--}5\times$ inference cost. The shared design has clear boundaries: the expensive options (geometry fusion, aggregation) pay off only on some applications, cluttered household objects need synthetic training data, heads do not transfer zero-shot, and deployment on real spacecraft images is currently limited by detection rather than by the heads. Fine-tuned state-of-the-art methods remain more accurate on every benchmark we use, because they adapt the backbone features to each target; our results show how much of that benefit can be recovered without any backbone gradient, at a small fraction of the adaptation cost and with backbones that stay shared. We report validation-selected results, matched comparisons and negative results, so that the cost and the limits of adapting frozen backbones with small heads can be judged directly.
